\documentclass[11pt]{article}
\usepackage[utf8]{inputenc}
\usepackage[T1]{fontenc}
\usepackage{lmodern}
\usepackage[margin=1in]{geometry}
\usepackage{amsmath,amssymb}
\usepackage{textcomp}
\usepackage[dvipsnames]{xcolor}
\usepackage{enumitem}
\usepackage{parskip}
\usepackage{titlesec}
\usepackage[colorlinks=true,linkcolor=blue,citecolor=blue,urlcolor=blue]{hyperref}
\usepackage[round]{natbib}
% U+2217 appears in one auto-generated refs.bib title (script-generated from the
% reference cache; entries must not be hand-edited) -- map it for pdflatex.
\DeclareUnicodeCharacter{2217}{\ensuremath{{}^{*}}}

\titleformat{\section}{\large\bfseries}{\thesection}{1em}{}
\titleformat{\subsection}{\normalsize\bfseries}{\thesubsection}{1em}{}
\titleformat{\subsubsection}{\normalsize\bfseries\itshape}{}{0em}{}

\title{\textbf{Checkpoints Certify a Signal, Not the State They Imply:\\ Measuring Bypass in Multi-Step Agent Verification}}
\author{}
\date{}

\newcommand{\ph}{\textit{[placeholder --- fill in after experiments are run]}}

\begin{document}
\maketitle
\vspace{-2em}

\begin{quote}
\small\textit{Assembled from the six confirmed stage outputs of an \texttt{auto-research-discussion} run (seed topic B1.3), post research-taste audit and post adversarial design review (\texttt{discussion/design\_review.md}, fixes applied 2026-07-04). Status: \textbf{specification / paper draft --- no experiments have been run.} All numbers below are pre-registered design targets, not results. Updated 2026-07-05 to fold in the design-review fixes: the estimand is split into a headline (natural-incidence) and secondary (constructibility-gate) quantity, the statistical plan models trajectory clustering within tasks, the natural-incidence power statement is corrected, the global-pivot machinery is re-targeted, and Related Work engages the three verified cache method-twins. Updated 2026-08-15 (see \texttt{discussion/review\_2026-08.md}): the motivating gap claim is repositioned against four Feb--Jun 2026 arXiv neighbors (verified via the arXiv export API) --- the generic ``verifier gaming is unmeasured'' framing is conceded as closed; the milestone-locality estimand and attribution claim is what remains and is claimed. RQ2's setup now mirrors the pre-registered completion-yield rule.}
\end{quote}

\begin{abstract}
Multi-step agent benchmarks and reward signals routinely grade long-horizon tasks by checking each intermediate milestone independently --- a file exists, a database row appears, a tool call returns the right signature --- because verifying full end-to-end correctness directly is expensive. Recent audits have shown that \emph{whole-task outcome} verifiers in agent benchmarks can be defeated at scale, and that unprompted agents already exploit verifiers in a measurable fraction of rollouts; what those audits leave unexamined is this practice's structural core at the \emph{milestone} level: a milestone verifier certifies a local, observable signal, but the prerequisite state that signal is supposed to imply is a property of causal history, and the map from history to signal is many-to-one. Many trajectories that never establish a milestone's true prerequisite state can still produce the same passing signal, so satisfying every checkpoint is strictly weaker than establishing what those checkpoints jointly certify. We build a multi-step agent task suite annotated with an explicit prerequisite dependency graph, equip each milestone with a verifier whose local-only scope is directly auditable in code, and score every trajectory against a separate terminal state-consistency oracle sharing no logic with the verifiers. Our headline estimand is the \emph{natural incidence} of checkpoint bypass in unprompted rollouts of real LLM agents; a constructed suite of out-of-order, step-skipping, and direct-injection trajectories serves as a constructibility gate and lets us characterize which strategies and graph topologies admit bypass at all, rather than being reported as a deployment incidence itself. A strongest-checker control isolates whether any observed bypass reflects structural locality rather than weak checker implementations, and a manual audit bounds the terminal oracle's own error before any rate is attributed to genuine bypass. Because trajectories are nested within tasks, every rate is reported with task-clustered (not trajectory-i.i.d.) intervals. We report the resulting checkpoint-bypass rates --- natural incidence, by strategy, and by graph topology --- and what they imply for milestone-based grading in agent evaluation and training.
\end{abstract}

\section{Introduction}

Long-horizon agent tasks are rarely graded end to end. A coding agent that must edit a file, run a migration, and update a config is instead checked milestone by milestone: does the file exist, did the test pass, does the database row match. This decomposition is not a shortcut taken reluctantly --- it is standard practice across agent benchmarks and training pipelines, because directly verifying whether the entire multi-step task actually succeeded is often expensive or underspecified, while each individual milestone typically has a cheap, well-defined local check; recent systems make the practice explicit, grading agents on verifiable milestone DAGs reconstructed from commit histories \citep{SWEMilestone2026} and anchoring RL training rewards to dynamically distilled milestones \citep{AdaptiveMilestoneReward2026}. The appeal is obvious: if every milestone's checker passes, the reasoning goes, the task must have been completed correctly.

That verifiers in agent settings can be gamed is, by mid-2026, no longer news: a recent audit found 16\% of 1,968 tasks across five terminal-agent benchmarks hackable by frontier models attacking the task's \emph{outcome verifier} \citep{HackerFixerLoops2026}, and an ultra-long-horizon benchmark observed unprompted reward-hacking behavior in 13.8\% of real-agent rollouts \citep{SWEMarathon2026}. But both of these measure the whole-task outcome check as a single attack surface. What still has not been directly measured is the reasoning that milestone decomposition quietly adds on top: that passing \emph{every} milestone checker is \emph{sufficient} evidence that the task's true prerequisite structure was actually respected. A milestone verifier can only observe its own local signal --- a file's presence, a row's content, an assertion's outcome --- and that signal is a many-to-one projection of the causal history that produced it. Many different histories, including ones that never executed the milestone's real prerequisite step, can produce an identical passing signal. A trajectory that completes milestones out of order, skips a step whose effect it fakes another way, or simply writes the expected file or database row directly, can in principle satisfy every local checker without ever establishing what those checkpoints are supposed to jointly certify. If this is possible in practice and not just in principle, then benchmark scores, training reward signals, and ``task complete'' judgments built on milestone-by-milestone grading may be systematically overstating real multi-step competence --- and the verifier itself, rather than the task, becomes what gets gamed.

The closest existing work approaches this concern from adjacent angles without closing it. The verifier-gaming measurement line just cited is nearest in spirit: it establishes that outcome verifiers are exploitable at scale \citep{HackerFixerLoops2026}, that exploitation can be made deterministically measurable by embedding hackable opportunities into environments by design \citep{HackVerifiableEnvs2026}, and that it occurs unprompted \citep{SWEMarathon2026} --- but in every case the object attacked is a single per-task outcome check, so none of these can ask whether a \emph{set} of milestone-scoped checkers over an explicit prerequisite graph can all pass while the end-to-end state is inconsistent, nor attribute any such gap to checker locality rather than checker weakness. Evaluation frameworks that catalog how agent benchmarks actually grade tasks --- file and existence checks, tool-signature checks, database-state comparisons, exact-match assertions --- argue convincingly that this local grading under-measures the interaction \emph{process}, but their remedy is a more holistic judge over the full trajectory, not a measurement of how often local checks can be satisfied while the true prerequisite state never holds. Work that builds agents around explicit prerequisite dependency graphs treats the graph as a planning aid for producing valid action sequences, correcting it from feedback when a plan fails --- but never asks the dual question of whether an independent milestone verifier can be satisfied while the graph's prerequisites are unmet. Elsewhere, process reward models that score reasoning steps independently are already known to be locally plausible yet globally inconsistent --- a step that looks fine on its own can become invalid once the full solution is assembled --- but this observation has so far been made for linear reasoning chains graded against a single final answer, not for multi-milestone agent trajectories with an explicit dependency structure and adversarial completion strategies. And checklist-based process reward models for web agents have shown that template-matching checkers can be satisfied by superficial compatibility rather than genuine task progress --- direct evidence that surface-level checking is exploitable, but not a measurement of how often, or under which conditions, that exploitation actually produces a checkpoint pass alongside true task failure. (This paragraph surveys works absent from our verified reference cache as of this draft --- see the Related Work cache-integrity note --- and is retained as anticipated positioning, not a supported citation.)

The nearest \emph{verified} neighbors --- method-twins with resolvable digests in our reference cache --- approach the same local-signal-versus-disjoint-oracle asymmetry from other domains: naturally produced code that passes a functional-correctness check yet fails a disjoint reproducibility oracle \citep{CodeReproducibilityGap2025}; step-level verification of reasoning chains because final-answer correctness can mask invalid intermediate steps \citep{RevealWeakestLink2024}; and theorem-proving pipelines where a correct final answer can coexist with logically invalid steps, caught only by a disjoint prover oracle \citep{MATPLogicalFlaws2025}. None constructs an explicit prerequisite dependency graph, none adversarially builds reorder/skip/inject trajectories over multi-milestone agent state, and none runs a strongest-checker attribution control to separate a structural finding from an implementation artifact.

Turning this gap into something measurable surfaces its own difficulties. The space of ways a trajectory could reorder, skip, or inject its way past a set of milestones grows far faster than the number of milestones, so no finite constructed suite can exhaustively cover it --- any reported rate on that suite is necessarily a sample, and has to be reported as a lower bound rather than a population-level number. Locality itself has to be established as a real, auditable property of the verifier code, not merely assumed, or any measured gap between milestone success and true task success could just as easily be an accident of a specific implementation. Among the three bypass strategies, direct state injection is the sharpest test of this locality: it forges a milestone's local success signal without needing an executable path through the actual task logic, and distinguishing a state that was legitimately produced by executing prior steps from one whose surface signature was simply written in place requires authenticating causal history --- something a lightweight, milestone-scoped checker cannot do without becoming as expensive as full state tracing. And because any measured bypass could, in principle, be an artifact of this study's own checkers being weaker than what a careful benchmark designer would actually deploy, the finding is only informative if it survives being tested against the strongest realistic local checker available, not just a convenient toy one. Finally, the terminal ground truth used to detect a bypass is itself a nontrivial specification --- a buggy oracle would manufacture false inconsistencies that masquerade as bypasses, or a lenient one would hide real ones.

Our approach turns each of these obstacles into a specific design decision. Task generation is stratified across prerequisite-chain depth and branching factor and across all three bypass-strategy families, so that any reported bypass rate can be decomposed by exactly the structural axes that determine where bypass concentrates, rather than presented as one undifferentiated number. Every per-milestone verifier is built with a statically auditable scope restriction --- no code path reads another milestone's provenance --- so that locality is a checkable property of the instrument, not an assumption about it. A dedicated injection-strategy probe holds a milestone's surface signature fixed while varying only whether its true prerequisite was executed, directly testing whether state injection defeats local checking as sharply as its logic predicts. A strongest-local-checker control replaces every verifier with the richest realistic check per checker type, still scoped to its own milestone, to test whether any bypass survives against the best a real system would plausibly deploy --- separating a structural finding from an artifact of weak instrumentation. And the terminal state-consistency oracle is validated on its own terms: a manual audit of its verdicts on trajectories known to be honest bounds its false-inconsistency rate before any bypass claim is reported. In its evidence shape, this is deliberately a controlled measurement study that directly tests an assumed-sufficient claim --- the mold of close empirical examinations such as the calibration analysis of pre-trained language models \citep{PLMCalibrationCloseLook2023} --- rather than a proposal for a competing verification method.

We evaluate this design on multi-step tasks spanning coding and filesystem operations, database and tool-use state, and web-page state, comparing constructed adversarial trajectories against honest, topologically-valid completions and against the strongest-checker control, and separately measuring how often real LLM agents produce a milestone-complete-but-inconsistent rollout without any adversarial intent --- this natural-incidence question is the paper's headline estimand; the constructed trajectories answer a different, prior question (can bypass be built at all, and which strategies and topologies admit it) and are never pooled with the natural-incidence rate.

In summary, we contribute:
\begin{itemize}[leftmargin=1.4em]
\item A formalization of the asymmetry between what a per-milestone verifier can observe and what it is meant to certify: milestone checking is a many-to-one projection of causal history onto a local signal, so satisfying every checkpoint is strictly weaker than establishing the prerequisite state those checkpoints jointly imply.
\item A constructed multi-step agent task suite with explicit prerequisite dependency graphs, auditably locality-scoped per-milestone verifiers mirroring standard benchmark grading, and a terminal state-consistency oracle held disjoint from and validated independently of those verifiers.
\item A strongest-checker attribution control that isolates whether observed bypass reflects the structural fact of milestone-scoped checking, rather than merely weak or lazy checker implementations --- together with an oracle-validity audit that bounds the ground truth's own error before any bypass rate is reported.
\item A headline measurement of the \emph{natural} checkpoint-bypass incidence in unprompted rollouts of real LLM agents, together with a secondary constructibility characterization --- by bypass strategy and by dependency-graph topology --- of the researcher-constructed trajectory space, with every rate reported under task-clustered (not trajectory-independent) uncertainty, and what both imply for milestone-based grading in agent benchmarks and reward design.
\end{itemize}

\section{Problem Definition}

\noindent\textbf{Formal setup.} Consider a multi-step task $T = (\Sigma, V, E, \Phi, \Omega)$: a state space $\Sigma$; a milestone set $V = \{v_1, \dots, v_n\}$; a prerequisite dependency graph $E \subseteq V \times V$; a set of per-milestone verifiers $\Phi = \{\varphi_k\}$; and a terminal oracle $\Omega$. A trajectory $\tau = (a_1, \dots, a_m)$ is an action sequence inducing a state path with terminal state $\sigma_\tau$. Each verifier $\varphi_k : \mathrm{view}_k(\tau) \to \{0,1\}$ is restricted to milestone $v_k$'s own local signal --- it has no read access to any other milestone's provenance --- and the all-pass indicator is $A(\tau) := \bigwedge_k \varphi_k(\tau)$. The oracle $\Omega : \sigma_\tau \to \{0,1\}$ is a total function comparing the trajectory's terminal state to the task's canonical goal state, sharing no code or logic with any $\varphi_k$. A \textbf{bypass event} is $B(\tau) := A(\tau)=1 \wedge \Omega(\tau)=0$: every milestone verifier passes, yet the terminal state is inconsistent with the goal.

The \textbf{checkpoint-bypass rate} is defined conditionally, $\rho := P_{\tau \sim D}[\Omega(\tau)=0 \mid A(\tau)=1]$, over a pre-registered trajectory distribution $D$, and is instantiated as \textbf{two distinct, never-pooled quantities}: the \textbf{headline} natural incidence $\rho_{\mathrm{nat}}$ over the distribution $D_{\mathrm{nat}}$ of real-agent rollouts (the reference population is rollouts a working milestone-verification pipeline would score complete; this is the quantity that bears on whether deployed per-milestone verification can be trusted without a separate terminal check), and the \textbf{secondary} constructed-suite rate $\rho_c$ over the researcher-authored distribution $D_c$ of adversarially built traces. Because $D_c$ is designed by construction to drive $A(\tau)=1$ while $\Omega(\tau)=0$ for the injection family, $\rho_c$ is an existence/constructibility quantity, not an incidence estimate of any deployed population --- it characterizes \emph{which strategies and topologies admit bypass}, and is reported as a manipulation-check gate rather than a headline finding.

\noindent\textbf{Why this is hard.} The bypass space grows combinatorially in the number of milestones --- reorderings, skip-patterns, and injection points over a dependency graph --- so any finite constructed sample can only characterize where bypass concentrates along structural axes such as prerequisite-chain depth and node in-degree, not certify an exhaustive or worst-case rate. Locality itself is not free to assume: $\varphi_k$'s restriction to $\mathrm{view}_k(\tau)$ has to be an auditable property of the verifier's code, since a claim about the consequence of milestone-scoped checking is only meaningful if that scoping is genuinely enforced rather than incidentally present. Among bypass mechanisms, direct state injection poses a distinctive difficulty: distinguishing a state produced by legitimately executing $v_k$'s prerequisites from one whose surface signature was written in place requires authenticating causal history, which a verifier confined to $\mathrm{view}_k(\tau)$ cannot do without effectively becoming a full state tracer. A measured $\rho_c > 0$ is also not on its own evidence that milestone-scoped checking is structurally insufficient, rather than merely evidence that this study's specific verifier implementations were weak --- ruling out the latter requires testing against the strongest realistic, still milestone-scoped checker available. And $\Omega$ itself is a nontrivial specification: since it must independently reconstruct and compare against a canonical goal state, its own false-inconsistency rate on trajectories known to be honest has to be bounded before any $B(\tau)$ count is attributed to genuine bypass rather than oracle error. Finally, because constructed trajectories are nested within tasks rather than independently sampled, every test below has to model that clustering explicitly rather than treating trajectories as i.i.d.\ Bernoulli draws.

\section{Method Overview}

At a conceptual level, this is a single measurement pipeline with five cooperating parts, applied to both adversarially constructed trajectories and to natural rollouts of real agents. Multi-step tasks are generated with an explicit prerequisite dependency graph and stratified across topology and bypass-strategy families; each milestone receives a verifier whose local-only scope is a checkable code property rather than an assumption; a dedicated injection probe and a strongest-checker control jointly test whether any measured gap between milestone success and true task success is a genuine consequence of that locality --- and specifically of state injection --- rather than an artifact of weak instrumentation; and a terminal oracle, held disjoint from every verifier, is itself audited before its verdicts are trusted as ground truth.

The \textbf{topology-stratified bypass sampler} parameterizes task generation over prerequisite-chain depth and branching factor and constructs reorder, skip, and injection trajectories across these strata, so the resulting bypass rate can be decomposed by exactly the structural features that plausibly drive it, rather than reported as a single undifferentiated number. The \textbf{locality-scoped verifier suite} implements each milestone's checker with a statically auditable restriction to that milestone's own local signal, mirroring the checker types real agent benchmarks use, and is the instrument applied to both constructed trajectories and real-agent rollouts to produce the constructed and natural bypass rates respectively. The \textbf{injection provenance probe} isolates the state-injection strategy specifically, holding a milestone's surface signature fixed while varying whether its true prerequisite was executed, to test whether injection defeats locality more sharply than reordering or skipping. The \textbf{strongest-checker attribution control} reruns the same measurement with the richest realistic, still milestone-scoped, checker per type, to test whether any bypass survives against the best real systems would plausibly deploy. And the \textbf{disjoint oracle with audit} validates the terminal ground truth itself, via a manual review of its verdicts on trajectories already known to be honest, before any bypass count is reported as evidence.

\section{Methodology}

\subsection{Topology-Stratified Bypass Sampling}
\textbf{Challenge.} The bypass space --- reorderings, skip-patterns, and injection points over an arbitrary prerequisite graph --- grows far faster than the number of milestones, so no finite constructed suite can exhaustively cover it, and an unstructured convenience sample gives no way to characterize where bypass concentrates or bound the direction of undercoverage.

\textbf{Insight.} Exhaustive coverage is impossible, but a \emph{stratified}, not exhaustive, sample --- organized along the structural axes (chain depth, branching factor) and strategy families most likely to determine where bypass concentrates --- turns ``we cannot cover everything'' into ``we can characterize where it concentrates and report the rate as an explicit, structured lower bound.''

\textbf{Method.} Task generation is parameterized over prerequisite-chain depth and node in-degree, and bypass trajectories are constructed across all three strategy families (reorder, skip, inject) within each topology stratum. The resulting checkpoint-bypass rate is always reported alongside its stratification --- never as a single pooled figure --- and explicitly labeled as a lower bound over the sampled distribution.

\subsection{Locality-Scoped Verifier Suite}
\textbf{Challenge.} A milestone verifier's job is to certify that milestone's prerequisite state, but if its scope restriction to that milestone's own local signal is merely assumed rather than enforced, any measured gap between milestone success and true task success could be an accident of one particular implementation rather than a consequence of milestone-scoped checking in general.

\textbf{Insight.} Locality has to be built into the verifier as a statically checkable code property --- no path reads another milestone's provenance --- so that it can be audited directly, and so that the same instrument can be pointed at two different trajectory distributions (adversarially constructed traces and natural agent rollouts) to produce two independently interpretable readings from one measurement design.

\textbf{Method.} Each milestone's verifier is implemented as a sandboxed check restricted to that milestone's own local signal --- file/existence checks, pattern or regex matches, unit sub-tests, tool-signature checks, database-state comparisons, and exact-match assertions, mirroring the mechanisms real agent benchmarks use --- with the scope restriction confirmed by a code-dependency audit before any rate is computed. The same verifier suite and disjoint terminal oracle are applied first to adversarially constructed trajectories, yielding the constructed checkpoint-bypass rate, and separately to unprompted rollouts of real LLM agents scored complete by the verifiers, yielding the natural bypass incidence.

\subsection{Injection Provenance Probe}
\textbf{Challenge.} Among the three bypass strategies, state injection is the one hardest to rule out in principle: it can forge a milestone's exact local success signal without needing to approximate any real control-flow path through the task, and distinguishing a legitimately produced prerequisite state from a directly written surface signature is, for a lightweight checker, effectively as hard as full causal-history tracing.

\textbf{Insight.} If injection is genuinely the sharpest bypass mechanism, holding a milestone's surface signature exactly fixed while varying only whether the real prerequisite operation was executed should isolate that mechanism cleanly, separate from any confound with the trajectory's overall plausibility as an executable action sequence.

\textbf{Method.} Within the injection-strategy trace family, each milestone's expected local artifact (the expected file, pattern, or database row) is written directly without executing the underlying prerequisite operation, and the resulting per-milestone verifier verdicts and bypass rate are compared against the reordering and step-skipping strategy families constructed over the same tasks.

\subsection{Strongest-Checker Attribution Control}
\textbf{Challenge.} Any measured bypass is only informative about milestone-scoped checking in general if it is not simply an artifact of this study's own verifier implementations being weaker or lazier than what a careful benchmark designer would realistically deploy --- otherwise the finding reduces to ``these particular toy checkers were easy to fool,'' a much less interesting claim.

\textbf{Insight.} Replacing each verifier with the richest realistic check of its type --- while preserving milestone-scoping by construction, so locality itself is never relaxed --- is the only way to test whether bypass survives against the strongest plausible real-world instrumentation, isolating structural locality from implementation weakness.

\textbf{Method.} Each per-milestone verifier is replaced by the fullest realistic local check of its checker type (the deepest unit sub-test, the strictest exact-match, the fullest database-state comparison), matched in capability and compute cost to the standard suite, while remaining scoped to its own milestone with no added cross-milestone visibility. The injection-strategy bypass rate is recomputed against this control and compared to its value under the standard verifier suite.

\subsection{Disjoint Oracle with Audit}
\textbf{Challenge.} A terminal oracle capable of detecting when milestone-level success does not imply true task success requires its own global specification of task state, which can be as difficult to get right as the original task's success criterion --- a buggy oracle would manufacture false inconsistencies that masquerade as bypasses, while an overly lenient one would hide genuine ones.

\textbf{Insight.} An oracle's trustworthiness is itself measurable, not merely assumable: trajectories that are known, by construction, to have honestly executed every prerequisite give a case where the ``correct'' verdict is known independently of the oracle, making them the right instrument for bounding the oracle's own error before its verdicts are trusted as ground truth for anything else.

\textbf{Method.} The terminal oracle is implemented as a total function of final task state --- a filesystem hash-set, a database row-set, or deterministic program output --- sharing no code or logic with any per-milestone verifier, confirmed by a code-dependency audit. A stratified sample of the oracle's verdicts on honest, topologically-valid trajectories is manually reviewed to measure its false-inconsistency rate, which is reported as an explicit bound on measurement error before any bypass rate involving that oracle is presented as evidence.

\section{Experiments}

\subsection{Experimental Setup}

\textbf{Datasets.} A constructed prerequisite-DAG agent task suite is the primary evaluation object: a hand-authored pilot set (coding and filesystem tasks, roughly 5--12 milestones per task, graded by file hashes and deterministic program output) establishes a canonical, independently computable terminal state cheaply, and a procedurally generated scale-up set extends coverage to tool-use/database-state and web-action domains. Real-world task substrates whose grading practice this suite's verifiers are built to mirror --- SWE-bench-style repository tasks, tool-use benchmarks with database-state, tool-selection, and tool-signature grading \citep{MCPAgentBenchToolUse2026,RLVRAtlassianToolUseAgents2026}, and function-calling leaderboards --- are used to justify that the constructed verifiers reflect standard, not idealized, agent-benchmark grading. Web-action tasks in the style of existing web-agent benchmarks are included specifically to test whether the topology-bypass relationship generalizes to a domain whose per-milestone checks are page/DOM-state comparisons rather than file, database, or tool-signature checks.

\textbf{Baselines} (sanity-check and attribution comparators, not competitors to beat --- this is a measurement study, and no research question here is framed as ``our procedure beats baseline X''). An honest-trajectory control --- topologically valid completions that genuinely execute every prerequisite --- is expected to produce a near-zero bypass rate; it is the necessary check that the verifier suite and the oracle agree on legitimate completions, so that any positive rate on constructed adversarial trajectories is attributable to the bypass mechanism rather than to baseline disagreement between the two measurement instruments. A strongest-local-checker control, described above as its own methodology module, tests whether bypass survives against the best realistic instrumentation rather than only against this study's standard checkers.

\textbf{Metrics.} The checkpoint-bypass rate reported for two distinct, never-pooled quantities --- the natural incidence $\rho_{\mathrm{nat}}$ (headline) and the constructed-suite rate $\rho_c$ (constructibility gate) --- broken down by bypass strategy and by graph-topology stratum; the odds ratio of bypass on prerequisite-chain depth and node in-degree from a mixed-effects logistic regression; and the oracle's own false-inconsistency rate from manual audit. Per-checker-type breakdowns are reported descriptively only --- the design pre-registers no inferential claim about which checker types are vulnerable (A3 is a continuous reporting discipline, not a hypothesis), an acknowledged scope restriction.

\textbf{Dependence structure (analysis-plan note).} Constructed trajectories are nested within tasks (roughly 5--15 traces per task share the same DAG, verifier code, and oracle), so no test below treats trajectories as independent draws. Rate intervals (RQ1, RQ2) use a task-level cluster bootstrap (10,000 resamples of tasks with replacement, percentile interval); strategy contrasts (RQ4) and the strongest-checker attribution ratio (RQ5) use cluster-robust (task-level) standard errors, equivalent to a GEE binomial model with exchangeable working correlation; the topology model (RQ3) is a mixed-effects logistic regression with a task-level random intercept. Every rate is reported with the number of distinct tasks ($n_{\mathrm{tasks}}$) alongside the number of trajectories, since $n_{\mathrm{tasks}}$ is the effective sample size under clustering, not $n_{\mathrm{trajectories}}$.

\textbf{Research questions:}
\begin{enumerate}[leftmargin=1.6em]
\item RQ1 --- Can checkpoint bypass be constructed at all (a constructibility gate, not the headline finding)?
\item RQ2 --- Do real agents produce checkpoint-complete-but-inconsistent rollouts without adversarial intent? (\textbf{headline})
\item RQ3 --- Does dependency-graph topology modulate the bypass rate?
\item RQ4 --- Which bypass strategy is most effective at defeating milestone verification?
\item RQ5 --- Is any observed bypass attributable to structural locality rather than weak checker implementations?
\item RQ6 --- Is the terminal oracle itself valid ground truth?
\end{enumerate}

\subsubsection{RQ1: Can checkpoint bypass be constructed at all? (constructibility gate)}
\textbf{Role.} RQ1 is a manipulation check, not a headline claim: because the injection strategy is defined to write the exact surface signature the milestone verifiers check, a pass is expected by construction. A pass licenses RQ3--RQ5 (which strategies/topologies admit bypass, and whether it survives stronger checkers) and the existence reading of Claim 1 (``bypass is constructible''); it is never read as evidence about deployed incidence --- that question belongs to RQ2. A failure here would be the more informative and surprising outcome.

\textbf{Setup.} On the pilot DAG suite, construct at least 300 trajectories (20--40 tasks $\times$ reorder/skip/inject $\times$ roughly 5 variants each) and compute the constructed-suite bypass rate $\hat\rho_c$ as the fraction passing every milestone verifier while failing the terminal oracle. Task-level cluster bootstrap (10,000 resamples, percentile interval) in place of a trajectory-i.i.d.\ binomial/Wilson interval.

\textbf{Expected outcome.} $\hat\rho_c \geq 0.10$ with the cluster-bootstrap 95\% lower bound above 0.05 --- demonstrating bypass is constructible on this suite, gating the strategy/topology/attribution analyses below.

\textbf{Falsification condition.} Refuted if the cluster-bootstrap 95\% lower bound on $\hat\rho_c$ does not exceed 0.05 --- that would mean constructed bypass trajectories cannot reliably defeat milestone verification on this suite even by design, an unexpected and highly informative failure.

\textbf{Hypothesis (plain language).} Milestone checkpoints can, in practice and not merely in principle, all be satisfied while the task's true prerequisite structure is never established --- a constructibility result, not yet a claim about deployed agents.

\textbf{Metrics.} Constructed-suite bypass rate $\hat\rho_c$ with task-level cluster-bootstrap 95\% CI and $n_{\mathrm{tasks}}$; per-checker-type breakdown.

\textbf{Result.} \ph

\subsubsection{RQ2: Do real agents produce checkpoint-complete-but-inconsistent rollouts without adversarial intent? (headline)}
\textbf{Role.} This is the paper's headline hypothesis: the reference population is rollouts a working milestone-verification pipeline would already score complete, so $\rho_{\mathrm{nat}}$ bears directly on whether deployed per-milestone verification can be trusted without a separate terminal check --- unlike RQ1's researcher-authored constructed rate. The named nearest precedent is \citet{SWEMarathon2026}'s 13.8\% unprompted reward-hacking incidence; $\rho_{\mathrm{nat}}$ is not comparable to that number by construction --- theirs counts detected exploit \emph{attempts} over all rollouts, ours counts terminal-state inconsistency \emph{conditional on} the rollout being scored milestone-complete, net of the audited oracle error (RQ6) --- and this estimand difference, not the existence of unprompted gaming, is what RQ2 contributes.

\textbf{Setup.} Among rollouts of 3--5 real LLM agents (at least 3 seeds each, spanning at least 20 distinct tasks) on the scale-up suite that are scored complete by the milestone verifier suite, measure the fraction that fail the terminal oracle, net of the audited oracle false-inconsistency rate (RQ6). Task-level cluster-bootstrap 95\% CI on $\hat\rho_{\mathrm{nat}}$, and a cluster-robust two-proportion comparison against $\hat\rho_c$ from RQ1, over at least 200 verifier-complete rollouts. \emph{Completion-yield rule (pre-registered in the experiment-design protocol):} the $\geq$200 floor counts \emph{verifier-complete} rollouts, not rollouts attempted, so the total rollout budget scales as 200 divided by the observed milestone-completion yield of the chosen agents on the scale-up tasks --- the walking-skeleton/pilot measures that yield first, gate PG-3 blocks scoring until the 200-complete/$\geq$20-task floor is met, and if the measured yield implies an infeasible total budget the task mix is re-selected under PG-0 rather than the floor being lowered. No yield value is assumed in advance.

\textbf{Decision rule (power-corrected).} The hypothesis holds iff $\hat\rho_{\mathrm{nat}} \geq 0.01$ ($\geq 2$ bypass events at $n=200$) with the cluster-bootstrap 95\% lower bound above 0. The decision threshold (0.01) is deliberately set \emph{below} the pre-registered effect floor ($\rho_{\mathrm{nat}} \geq 0.02$, the smallest incidence considered deployment-relevant) because setting the threshold equal to the floor caps power at $\approx 0.57$ regardless of $n$; decoupling them gives power $P(X \geq 2 \mid \mathrm{Bin}(200, 0.02)) \approx 0.91$ at the floor. Pre-declared edge cases: $\hat\rho_{\mathrm{nat}} = 0 \to$ FAIL (invoke the RQ2 fallback); $0 < \hat\rho_{\mathrm{nat}} < 0.01 \to$ INCONCLUSIVE-at-$n$ (report the required $n$), never reported as a null result.

\textbf{Expected outcome.} Natural incidence clears the decision rule above and is significantly lower than the constructed-suite rate --- bypass occurs during ordinary operation, but less often than under deliberate construction.

\textbf{Falsification condition.} Refuted (FAIL) if $\hat\rho_{\mathrm{nat}} = 0$; INCONCLUSIVE-at-$n$ if $0 < \hat\rho_{\mathrm{nat}} < 0.01$. Either non-PASS outcome means bypass is not shown to already occur, unprompted, in ordinary agent operation at the pre-registered floor.

\textbf{Hypothesis (plain language).} Checkpoint bypass is not purely a constructed vulnerability; it already occurs, at a lower but nonzero rate, in agents that were never instructed to exploit the verification scheme.

\textbf{Metrics.} Natural bypass incidence $\hat\rho_{\mathrm{nat}}$ with task-level cluster-bootstrap 95\% CI and $n_{\mathrm{tasks}}$; cluster-robust two-proportion comparison against $\hat\rho_c$.

\textbf{Result.} \ph

\subsubsection{RQ3: Does dependency-graph topology modulate the bypass rate?}
\textbf{Setup.} Over the scale-up suite (at least 200 tasks), fit a \textbf{mixed-effects logistic regression} of the per-trajectory bypass outcome on prerequisite-chain depth and node in-degree, with a \textbf{task-level random intercept} (trajectories are nested in tasks; see the dependence-structure note above). Wald tests on the fixed-effect coefficients, Holm-corrected across the two predictors, $\alpha = 0.05$.

\textbf{Expected outcome.} At least one fixed-effect coefficient shows an odds ratio greater than 1.2 per unit, with a significant Wald test --- deeper or higher-fan-in prerequisite structure raises the odds of bypass.

\textbf{Falsification condition.} Refuted if neither predictor's odds ratio exceeds 1.2 with a Holm-corrected significant Wald test --- that would mean bypass risk is not modulated by prerequisite-chain depth or node in-degree, i.e., graph topology does not predict where bypass concentrates.

\textbf{Hypothesis (plain language).} The risk of checkpoint bypass is not uniform across task structures; it is amplified by more complex prerequisite graphs, which give a locality-scoped verifier more unseen prior state to be decoupled from.

\textbf{Metrics.} Odds ratio per unit depth and per unit in-degree (mixed-effects logistic regression, task random intercept), with confidence intervals and Wald test statistics; $n_{\mathrm{tasks}}$ and per-cell trajectory counts.

\textbf{Result.} \ph

\subsubsection{RQ4: Which bypass strategy is most effective at defeating milestone verification?}
\textbf{Setup.} Compare per-strategy bypass rates (reorder, skip, inject) via three pairwise two-proportion tests using \textbf{cluster-robust (task-level) standard errors} (equivalently, GEE binomial with strategy as predictor and exchangeable working correlation), Holm-corrected across the three comparisons, $\alpha = 0.05$.

\textbf{Expected outcome.} The injection strategy's bypass rate exceeds both reordering and step-skipping by at least 10 percentage points.

\textbf{Falsification condition.} Refuted if the injection strategy's bypass rate does not exceed both reordering's and step-skipping's rates by at least 10 percentage points (Holm-corrected, cluster-robust pairwise tests) --- that would mean forging a milestone's surface signature directly is not a distinctively stronger bypass than merely reordering or omitting steps.

\textbf{Hypothesis (plain language).} Directly forging a milestone's local success signal is a more effective bypass than merely reordering or omitting steps, because injection is not constrained by needing an executable path through the task's real control flow.

\textbf{Metrics.} Per-strategy bypass rate with cluster-robust pairwise contrasts, Holm-corrected p-values, and $n_{\mathrm{tasks}}$ per strategy cell.

\textbf{Result.} \ph

\subsubsection{RQ5: Is any observed bypass attributable to structural locality rather than weak checker implementations?}
\textbf{Setup.} Recompute the injection-strategy bypass rate under the strongest-local-checker control (the richest realistic, still milestone-scoped check per checker type) and compare it to the rate under the standard verifier suite, using a one-sided two-proportion non-inferiority test with \textbf{cluster-robust (task-level) standard errors} --- the strong- and weak-checker rates are computed on the same tasks, so pairing by task is mandatory, not optional.

\textbf{Expected outcome.} The bypass rate under the strongest-checker control retains at least half its value under the standard suite (cluster-robust lower bound clears 0.5) --- the bypass is not eliminated by using better, still milestone-scoped, checkers.

\textbf{Falsification condition.} Refuted if the cluster-robust lower bound on the ratio (strongest-checker rate / standard-suite rate) falls below 0.5 --- that would mean the measured bypass is substantially attributable to weak or lazy checker implementations rather than to the structural fact of milestone-scoped checking.

\textbf{Hypothesis (plain language).} The vulnerability being measured is a consequence of checking milestones in isolation at all, not a symptom of this study's checkers being unusually weak.

\textbf{Metrics.} Bypass-rate ratio (strongest-checker control / standard suite) with a cluster-robust (task-paired) confidence interval.

\textbf{Result.} \ph

\subsubsection{RQ6: Is the terminal oracle itself valid ground truth?}
\textbf{Setup.} Manually audit a stratified sample of 50--100 oracle verdicts on honest, topologically-valid trajectories (which should uniformly pass both the verifier suite and the oracle), and separately confirm the honest-trajectory bypass rate is near zero.

\textbf{Expected outcome.} The oracle's false-inconsistency rate on honest trajectories is at most 5\%, and the honest-trajectory bypass rate is below 1\%.

\textbf{Falsification condition.} Refuted if the manual audit finds a false-inconsistency rate above 5\% on honest trajectories, or if the honest-trajectory bypass rate is at or above 1\% --- either would mean the terminal oracle is not yet valid ground truth, and any bypass rate measured against it cannot be attributed to genuine paradigm insufficiency rather than oracle noise.

\textbf{Hypothesis (plain language).} The terminal oracle is a trustworthy, independently validated ground truth, not a second unaudited measurement whose own errors could be mistaken for genuine bypasses.

\textbf{Metrics.} Oracle false-inconsistency rate from manual audit; honest-trajectory bypass rate.

\textbf{Result.} \ph

\vspace{0.5em}
\noindent\textbf{Further evaluation.} With more time, a natural extension is to test a holistic, non-milestone-scoped verification design --- a single judge scoring the whole trajectory at once, in the spirit of process-supervision approaches built for linear reasoning chains --- as an alternative arm alongside the milestone-scoped suite, to ask whether abandoning locality altogether (at a likely cost in scalability and per-benchmark portability) is the only way to close the measured gap, or whether a cheaper partial-provenance check suffices. This is deliberately out of scope here because it changes the object of study from measuring an existing paradigm's sufficiency to designing and validating a competing one, which would need its own outcome-tier evaluation rather than an extension of the measurement instrument above.

\section{Related Work}

The closest prior work to this paper spans three groups: a \textbf{verifier-gaming measurement group} of 2026 arXiv works (verified by ID against the arXiv API, 2026-08-15) that are now the nearest neighbors overall and that this paper concedes to and positions against first; a \textbf{verified-cache group} of three method-twins that are \texttt{selected} cache records with resolvable digests; and a \textbf{disjoint-corpus group} (TED, XENON, Gistify, StepORLM, LatentVeracity, WebArbiter, GTool, DiffTaskGraph) surfaced by an early keyword search but \textbf{absent from this topic's verified reference cache} --- retained below only as named future harvest targets, per the cache-integrity note logged in \texttt{digest.md} and \texttt{discussion/pipeline.md}, and not load-bearing for any claim in this paper.

\textbf{Verifier-gaming measured at scale (2026, arXiv-verified).} \citet{HackerFixerLoops2026} is the nearest neighbor to this paper's motivating claim: it audits 1,968 tasks across five terminal-agent benchmarks, finds 16\% hackable by frontier models attacking the task's outcome verifier, and proposes an adversarial hacker--fixer loop that patches verifiers until known exploits fail. We concede outright that verifier exploitability in agent benchmarks is therefore already directly measured --- the gap this paper claims is strictly narrower and different in kind: their unit is a single whole-task outcome verifier and their estimand is attack success given the task description, whereas ours is the conditional rate at which \emph{every} milestone-scoped verifier over an explicit prerequisite DAG passes while a disjoint, independently audited terminal-state oracle fails, decomposed by bypass strategy and graph topology, with a strongest-checker control attributing any gap to checker \emph{locality} rather than checker \emph{weakness} --- an attribution question their fixer loop's own success (driving attack rates to zero by patching implementations) sharpens but does not answer. \citet{SWEMarathon2026} supplies the natural-incidence precedent: 13.8\% of unprompted real-agent rollouts on ultra-long-horizon software tasks exhibit reward-hacking behavior detected by adversarial review of test suites and environments. We concede that unprompted verifier exploitation is thereby established as a live phenomenon; our headline $\rho_{\mathrm{nat}}$ differs operationally --- it is conditioned on rollouts a milestone-verification pipeline would already score complete, scored against a terminal state-consistency oracle rather than by detecting exploit \emph{attempts}, and reported net of an audited oracle false-inconsistency rate --- so the two numbers answer different questions (how often agents try to game a grader vs.\ how often milestone-complete rollouts are terminally inconsistent). \citet{HackVerifiableEnvs2026} pioneers the constructed-arm paradigm this paper's RQ1 also uses --- embedding detectable exploitation opportunities into environments so exploitation is verifiable by design --- in TextArena game environments with a single evaluation signal per environment; our constructed suite transplants that design move to multi-milestone, prerequisite-structured agent tasks where the object of study is the checker \emph{set}, not one signal. Finally, \citet{SWEMilestone2026} and \citet{AdaptiveMilestoneReward2026} document that milestone-DAG grading and milestone-anchored RL rewards are now deployed practice --- they use milestone verification as an instrument and assume its sufficiency; this paper audits exactly that assumption. None of the five constructs locality-audited per-milestone verifiers, none defines a conditional bypass estimand against a disjoint validated terminal oracle, and none runs a locality-vs-weakness attribution control.

\textbf{Verified method-twins (cache-grounded).} \citet{CodeReproducibilityGap2025} is the nearest verified structural neighbor: it shows AI-generated code that passes a functional-correctness benchmark --- the assumed-sufficient local signal --- still fails a separate, disjoint reproducibility oracle built from clean-environment re-execution. We concede the shared measurement pattern (a local signal passes while a disjoint oracle fails) outright; our contribution is an explicit prerequisite DAG over milestones, \emph{adversarially constructed} reorder/skip/inject trajectories rather than naturally produced code, milestone granularity rather than whole-program granularity, and a strongest-checker attribution control (RQ5) that isolates structural locality from implementation weakness. \citet{RevealWeakestLink2024} (REVEAL) introduces step-level verification of chain-of-thought reasoning specifically because final-answer correctness can mask invalid intermediate steps --- the same local-versus-global verification-granularity gap this paper studies, but in the \emph{inverse} direction (global pass masking local invalidity, rather than local pass masking global invalidity), over linear text chains rather than stateful agent trajectories with a prerequisite structure, and without a measured bypass rate under controlled adversarial strategies. \citet{MATPLogicalFlaws2025} measures a comparable gap in automated theorem proving --- a correct final answer can coexist with logically invalid steps --- using a disjoint theorem-prover oracle; we concede the disjoint-oracle validation pattern, and differ in operating on agent state (filesystem, database, tool-call state) rather than proof text, in treating prerequisite-state \emph{causality} (was the state legitimately produced, i.e., provenance) as the object of study rather than logical validity of a static artifact, and in decomposing the measured rate by bypass strategy and graph topology. \citet{RewardModelSelectionCrisis2025} shows a Goodhart-style gap in personalized alignment: standard reward-model ranking accuracy --- a locally-assumed-sufficient proxy --- fails to predict true downstream behavioral alignment under reward-guided decoding. We concede this is the same assumed-sufficient-checker-vs-true-outcome asymmetry, in single-turn decoding rather than multi-step agent trajectories; our contribution is the explicit prerequisite DAG, the multi-milestone dependency structure, and the adversarially constructed (rather than naturally sampled) bypass traces. None of the four constructs an explicit prerequisite dependency graph, none adversarially constructs reorder/skip/inject trajectories, and none runs a strongest-checker attribution control --- the combination this paper's design contributes.

\textbf{Agent-benchmark grading mechanisms and their remedies (disjoint-corpus, unverified --- named harvest targets only).} Work cataloging how LLM-agent benchmarks are actually graded documents that evaluation is fragmented across environment-specific mechanisms --- database-state comparisons, tool-signature checks, exact-match assertions --- that are reliable within a single benchmark but reduce rich multi-turn interaction to a final pass/fail signal, and proposes replacing them with a holistic LLM judge scoring natural-language grading notes over the full trajectory. Checklist-based process reward models for web agents make a complementary point: template-matching checkers can be satisfied by superficial compatibility rather than genuine progress, motivating a reasoning-first, pairwise-preference alternative. And execution-faithful benchmarking work shows that repository-level coding evaluations can be passed via memorized associations or retrieval shortcuts rather than real capability, motivating a benchmark tied to observed runtime behavior instead of static artifacts. All three name real weaknesses of milestone-style or template-style checking and propose a better \emph{checker} or grading protocol --- none isolates and measures the \emph{rate} at which currently standard, milestone-scoped checkers can all pass while a separately computed terminal ground truth fails. (These works could not be located in the verified reference cache as of this draft; the positioning above is retained as an anticipated harvest target, not a supported citation, and no claim in this paper depends on it.)

\textbf{Prerequisite-dependency planning (disjoint-corpus, unverified --- named harvest targets only).} Agents built around explicit prerequisite dependency graphs --- one for open-world planning, correcting an item-dependency graph from execution feedback, and one for tool selection, injecting a request-conditioned tool-dependency graph into a frozen model --- both show that violating prerequisite structure produces plans that are invalid despite looking locally plausible, and that errors in an early prerequisite invalidate every descendant step. Both treat the dependency graph as a planning aid for producing or correcting valid action sequences; neither audits the dual question this paper asks: whether an independent, milestone-scoped verifier can be satisfied while the graph's prerequisites are unmet in the first place. A third line of work learns a prerequisite adjacency structure directly from observed procedural video to detect ordering mistakes online --- a candidate ingredient for building a prerequisite-aware checker, but one whose detection is itself a learned forecasting signal that presumes the learned graph is correct, rather than an independent oracle over ground-truth task state. (Same cache-verification caveat as above.)

\textbf{Step-wise and process-level verifier locality (disjoint-corpus, unverified --- named harvest targets only).} Process reward models that score reasoning steps independently are already documented to be myopic: a step that looks locally correct can become inconsistent once the full solution is assembled, since correctness can depend on later context the step-wise scorer never sees. A related line of work on latent veracity inference makes the same granularity point from the other direction --- answer-conditioned verifiers score a whole candidate solution rather than localizing failure to individual steps. Both identify essentially the same local-versus-global tension this paper measures, but for linear reasoning or derivation chains scored against a single final answer, with no explicit prerequisite dependency graph, no adversarially constructed bypass trajectory, and no independent state-consistency oracle held disjoint from the scoring mechanism. (Same cache-verification caveat as above --- retained as a harvest target, not a supported citation.)

\section{Conclusion}

This work turns an assumption baked into standard agent-benchmark grading --- that passing every milestone checkpoint is good evidence a multi-step task actually succeeded --- into a measurable question with a precise identifiability structure, rather than treating milestone-based verification as self-evidently adequate. Whole-task outcome-verifier gaming is already measured at scale \citep{HackerFixerLoops2026,SWEMarathon2026}; what this design adds is the milestone-level version of that question --- whether an entire set of locality-scoped checkers over an explicit prerequisite DAG can pass while an independently validated terminal state fails, and whether that is attributable to locality itself. The central methodological wager is that a checkpoint-bypass rate is only informative once locality is built into the verifiers as an auditable property, the strongest realistic checker has been tried and still bypassed, and the terminal ground truth used to detect the gap has itself been validated against trajectories known to be honest; a weaker design risks attributing a measured rate to the wrong cause, or missing it entirely behind oracle noise. The paper's headline claim is about the \emph{natural} incidence $\rho_{\mathrm{nat}}$ in unprompted real-agent rollouts (RQ2), not the researcher-constructed rate $\rho_c$ (RQ1), which is a manipulation-check gate expected to pass by construction; the two are never pooled. Whichever way $\rho_{\mathrm{nat}}$ and the strongest-checker attribution (RQ5) land, the finding is decision-relevant either way. If natural incidence clears its pre-registered floor and survives the strongest-checker control, benchmark scores, training rewards, and deployment ``task complete'' judgments built on milestone checking need a companion terminal-consistency mechanism before they can be trusted at face value. If natural incidence is null or the bypass is substantially attributable to weak checker implementations (RQ5 fails), the pre-committed downgrade applies: checkpoint bypass is a constructible artifact of milestone-locality that current real-agent behavior does not exhibit unprompted and that stronger, still milestone-scoped, checkers substantially mitigate --- a design caution rather than a measured deployment risk, per the paper's pre-registered global-pivot wording. Three limitations bound this claim from the outset: the constructed-suite rate is inherently a sampled lower bound over a combinatorially large bypass space, not a worst-case or population-level figure; every reported rate is computed under task-level clustering (trajectories are nested in tasks, not i.i.d.), so uncertainty intervals are wider and effective sample sizes smaller than a naive per-trajectory count would suggest; and the study is scoped to domains with a canonical, independently computable terminal state, leaving open-ended tasks without such a state for future work on oracle design.

\vspace{1em}
\hrule
\vspace{1em}

\section*{Pre-handoff checklist}
\begin{itemize}[leftmargin=1.4em]
\item \textbf{Organization/navigation} --- narrative flows problem $\to$ method $\to$ experiments $\to$ related work without forward references to undefined terms.
\item \textbf{Terminology discipline} --- ``checkpoint-bypass rate,'' ``locality-scoped verifier,'' ``terminal oracle,'' ``prerequisite state'' each defined at first use in Problem Definition, used consistently thereafter; no internal stage IDs (\texttt{sp\_*}, \texttt{m\_*}, \texttt{rq\_*}) leaked into prose --- checked directly against the discussion pipeline log.
\item \textbf{Notation with worked example} --- the formal setup section defines $(T, \Sigma, V, E, \Phi, \Omega, \tau, \sigma_\tau, A(\tau), B(\tau), \rho_c, \rho_{\mathrm{nat}})$ and the Problem Definition/Methodology sections consistently reuse this notation; no separate worked numeric example is included since no numbers exist yet (would require fabricating results).
\item \textbf{RQ framing matches outcomes} --- all 6 RQs are graded/quantitative (rates, odds ratios, agreement/retention proportions), each answered in its own subsection with an explicit falsifiable expected outcome, not a binary yes/no.
\item \textbf{No typesetting artifacts} --- only explicit placeholder markers in Results, per the drafting convention this is a pre-registered spec, not a report.
\item \textbf{Artifact plan} --- the constructed task suite (dependency graphs, verifiers, oracles) and the audited real-agent rollout transcripts are the natural release artifacts; not yet built, so not claimed as available.
\item \textbf{Bibliography wiring} --- \texttt{refs.bib} is generated from \texttt{discussion/reference\_candidates.jsonl}'s \texttt{selected} records (17 entries after the 2026-08-15 review pass appended six arXiv-verified neighbors); every \texttt{selected} \texttt{method-twin} record is cited with a real \texttt{\textbackslash citep}/\texttt{\textbackslash citet} in Related Work, not paraphrased without attribution.
\end{itemize}

\section*{Hand-off}
This is a draft, not a finished paper --- experiments have not run. Natural next steps: prose polish and venue-specific LaTeX formatting once a target venue is chosen, or revising an earlier stage if piloting surfaces a problem (e.g., if the honest-trajectory bypass rate comes in above the pre-registered 1\% gate, an oracle-validity failure).

\bibliographystyle{plainnat}
\bibliography{refs}

\end{document}
